\scanpage{11}
\uline{Corollaire} 1 \quad Il y a correspondance biunivoque entre les idéaux $J$ de $\overset n{\bigotimes}_A B$ contenus dans l'idéal d'augmentation, et les idéaux $J^{\unsure{0|o}}$ de $\overset n{\bigotimes}_{A_0}B_0$ contenus dans l'\unsure{idéale|idéal} d'augmentation, donnée par
\[
\left\{\begin{aligned}
J&\longmapsto\unsure{\rho_n|p_n}(J)\\
\rho_n^{-1}(J^0)\cap\Delta_{B/A}^{(n)}&\longleftarrow J^0
\end{aligned}\right.
\]
Pour deux idéaux \inserted{$J,J^0$} qui se correspondent, l'hom naturel $J\longrightarrow J^{\unsure{0|2}}$ est bijectif.

Noter que la correspondance précédente est compatible avec les relations d'ordre, \deleted{et \unsure{la}} \inserted{\unsure{pour|puis}} avec la formation de $+$ et $-$. D'autre part, les \deleted{\unsure{formations de ces}} \inserted{\unsure{applications|bijections|identifications} précédentes} (pour $n$ variable) commutent aux opérations simpliciales, dans un sens évident. \inserted{(cf Prop 1, cor 1)}

On conclut de \deleted{\unsure{ceci}} \inserted{ces remarques}

\uline{Corollaire} 2 \quad \unsure{(a)|($\alpha$)} Il y a correspondance biunivoque (induite par celle du cor. 1 pour $n=2$) entre relations d'équivalence dans
% archive sheet number: 10
